\chapter{Group B Streptococcal Infection}

\section*{Definition and Overview}
Group B streptococcus (GBS), also called \textit{Streptococcus agalactiae}, is a bacterium that commonly lives harmlessly in the bowel, vagina, and urinary tract of healthy adults. It is best known for its effect on pregnancy: around 1 in 5 pregnant women carry GBS, and although it rarely causes problems in the mother, it can be passed to the baby during birth and cause serious illness in newborns, including pneumonia, sepsis, and meningitis. GBS can also cause infection in older adults and people with weakened immunity. Pregnant women are screened for GBS late in pregnancy, and those who carry it are offered antibiotics during labour to protect the baby. GBS infection in newborns is serious but largely preventable.

\section*{Epidemiology}
Around 15--25\% of pregnant women carry GBS in the vagina or bowel at the time of delivery. Without preventive treatment, around 1 in 200 babies born to carriers develop early-onset GBS disease within the first week of life. With intrapartum antibiotics, the risk falls to about 1 in 4,000. In Australia, GBS is a leading cause of serious infection in newborns, affecting roughly 1 in 1,000 live births, and the risk is higher in premature babies. GBS also causes infection in older adults and people with diabetes or other chronic illness, with rates rising in recent decades.

\section*{Aetiology and Pathophysiology}
GBS causes disease by colonising the mother and being transmitted to the baby around birth.

\begin{itemize}
    \item \textbf{Colonisation:} GBS lives harmlessly in the bowel and genital tract of many healthy women
    \item \textbf{Transmission to baby:} the bacterium can be passed to the baby during vaginal delivery, mainly after the membranes rupture
    \item \textbf{Early-onset disease:} infection in the first week of life, usually within 24--48 hours, causing pneumonia, sepsis, or meningitis
    \item \textbf{Late-onset disease:} infection from 7 days to 3 months, often meningitis, sometimes transmitted from sources other than the mother
    \item \textbf{Risk factors:} carrying GBS, premature birth, prolonged rupture of membranes, fever in labour, and a previous baby with GBS disease
    \item \textbf{Disease in adults:} GBS can infect the skin, joints, urinary tract, and bloodstream, especially in older people and those with chronic disease
    \item \textbf{Prevention:} intrapartum antibiotics given to carriers during labour interrupt transmission to the baby
\end{itemize}

\section*{Clinical Features}
Symptoms depend on the age and health of the affected person.

\begin{itemize}
    \item In newborns (early-onset): poor feeding, lethargy, irritability, grunting, fast or laboured breathing, fever or low temperature, and a blotchy or grey appearance
    \item In newborns (late-onset): fever, poor feeding, irritability, and signs of meningitis such as a high-pitched cry and a bulging fontanelle
    \item In pregnant women: usually no symptoms; occasionally urinary tract infection or infection of the womb after birth
    \item In adults: urinary tract infections, skin and soft tissue infections, joint infections, and rarely bloodstream infection
    \item Serious disease in adults: fever, chills, confusion, and features of sepsis
\end{itemize}

\section*{Differential Diagnosis}
Other causes of newborn and adult infection should be considered.

\begin{itemize}
    \item Other bacterial causes of newborn sepsis, including \textit{E. coli}, \textit{Listeria}, and \textit{Staphylococcus aureus}
    \item Viral infections in newborns, including herpes simplex virus
    \item Other causes of neonatal respiratory distress, including respiratory distress syndrome and meconium aspiration
    \item In adults: other causes of urinary tract, skin, and bloodstream infection
    \item Meningitis from other bacteria
    \item Non-infectious causes of fever and poor feeding in newborns
\end{itemize}

\section*{When to Seek Medical Care}
Pregnant women should attend antenatal care, including GBS screening in late pregnancy, and discuss their GBS status with their maternity team. A newborn who is feeding poorly, lethargic, or breathing abnormally needs urgent medical review.

\textbf{Call 000 (or local emergency services) for:}
\begin{itemize}
    \item A newborn with fever, poor feeding, lethargy, difficulty breathing, or seizures, which may indicate serious infection and requires emergency care
    \item Signs of meningitis in a baby, including a high-pitched cry, stiff body, or bulging soft spot
    \item Fever in labour or rupture of membranes with signs of infection
    \item Any adult with fever, confusion, or rapid deterioration
\end{itemize}

\section*{Diagnosis and Investigations}
Diagnosis relies on microbiology tests.

\begin{itemize}
    \item \textbf{Prenatal screening:} a vaginal and rectal swab at 35--37 weeks of pregnancy detects GBS carriage
    \item \textbf{Blood cultures:} the key test in suspected newborn or adult invasive disease
    \item \textbf{Lumbar puncture:} examines cerebrospinal fluid when meningitis is suspected
    \item \textbf{Chest X-ray:} for pneumonia
    \item \textbf{Urine tests:} for urinary tract infection
    \item \textbf{Full blood count and inflammatory markers:} support the diagnosis of infection
\end{itemize}

\section*{Management}
The mainstay of management is prevention during labour, with treatment of established disease.

\begin{itemize}
    \item \textbf{Intrapartum antibiotics:} penicillin given intravenously during labour to GBS carriers prevents transmission to the baby
    \item \textbf{Screening-based approach:} all women are screened at 35--37 weeks, and carriers receive intrapartum antibiotics
    \item \textbf{Risk-based approach:} in some settings, antibiotics are given when risk factors are present rather than through universal screening
    \item \textbf{Treatment of newborn disease:} hospital care with intravenous antibiotics, usually penicillin and gentamicin, plus supportive care
    \item \textbf{Treatment of adult disease:} antibiotics appropriate to the site of infection
    \item \textbf{Monitoring:} babies at risk are observed for signs of infection after birth
    \item \textbf{Follow-up:} babies treated for meningitis need hearing and developmental follow-up
\end{itemize}

\section*{Complications}
\begin{itemize}
    \item Early-onset neonatal disease: pneumonia, sepsis, and meningitis, which can be fatal or cause lasting disability including hearing loss and developmental delay
    \item Late-onset neonatal meningitis and its complications
    \item Maternal infections after birth, including endometritis and urinary tract infection
    \item Chorioamnionitis, infection of the membranes during labour
    \item Invasive disease in older adults, including sepsis and joint infections
    \item Adverse effects of antibiotics, including maternal allergic reactions (rare)
\end{itemize}

\section*{Prognosis and Outcomes}
The outlook for GBS infection has improved dramatically with preventive treatment. Babies born to screened and treated mothers have a very low risk of GBS disease, and when early-onset disease does occur, prompt intensive care is highly effective, with most babies recovering fully. The mortality of early-onset GBS disease has fallen to around 3--5\%, though meningitis can still cause long-term disability in a minority. Screening and intrapartum antibiotics are among the most successful preventive interventions in obstetrics.

\section*{Prevention and Self-Care}
\begin{itemize}
    \item Attend antenatal care, including GBS screening at 35--37 weeks of pregnancy
    \item Discuss your GBS status and birth plan with your maternity team
    \item If you are a carrier, ensure your labour and birth plan includes intrapartum antibiotics
    \item Report fever, leaking membranes, or signs of infection in labour promptly
    \item Know the signs of newborn illness: poor feeding, lethargy, fever, and breathing problems
    \item Seek urgent medical care if you are worried about your baby at any time
    \item Older adults and people with chronic illness should seek prompt treatment of infections and review of fevers
\end{itemize}

\section*{Sources and Further Reading}
The following reputable sources provide further information for healthcare students and patients seeking reliable, current advice about group B streptococcal infection.

\begin{itemize}
    \item Royal Australian and New Zealand College of Obstetricians and Gynaecologists. \textit{GBS screening and prevention in pregnancy.} https://ranzcog.edu.au/
    \item Healthdirect Australia. \textit{Group B streptococcus in pregnancy.} https://www.healthdirect.gov.au/
    \item Pregnancy, Birth and Baby (Australian Government). \textit{GBS and pregnancy.} https://www.pregnancybirthbaby.org.au/
    \item Australian Government Department of Health and Aged Care. \textit{Group B streptococcal disease fact sheet.} https://www.health.gov.au/
    \item NHS. \textit{Group B strep.} https://www.nhs.uk/
    \item Centers for Disease Control and Prevention. \textit{Group B strep.} https://www.cdc.gov/
\end{itemize}
